% =====================================================================
\chapter{Integrales impropias}
\label{tema:4}
% =====================================================================

La integral de Riemann exige un intervalo cerrado y acotado y una función acotada. Muchas integrales importantes no cumplen alguna de las dos condiciones: la densidad normal está definida en toda la recta real, la función Gamma integra en $[0,\infty)$ y la densidad de una ley de potencia no está acotada cerca del origen si el exponente es negativo. Este tema extiende la integral mediante un paso al límite, estudia cuándo ese límite existe y calcula dos integrales fundamentales para la probabilidad: la función Gamma y la integral gaussiana. Se siguen la §3.7 de~\cite{openstax2}, la §7.5 de Strang~\cite{strang}, el capítulo~10 de Apostol~\cite{apostol} y las sesiones~91 a~93 del curso 18.01SC~\cite{mit1801}.

\section{Integrales en intervalos no acotados}

\begin{definicion}[Integral impropia de primera especie]\label{def:impropia1}
Sea $f\colon[a,\infty)\to\R$ integrable en $[a,b]$ para todo $b>a$. Si existe el límite finito
\[
\int_a^{\infty}f(x)\dif x=\lim_{b\to\infty}\int_a^{b}f(x)\dif x,
\]
se dice que la integral impropia \emph{converge} y que ese es su valor; en otro caso, \emph{diverge}. Se define igual $\int_{-\infty}^{b}f$. Para una función definida en todo $\R$, se dice que $\int_{-\infty}^{\infty}f$ converge si convergen $\int_{-\infty}^{c}f$ y $\int_{c}^{\infty}f$ para algún $c$, y su valor es la suma de ambas.
\end{definicion}

Por la aditividad, la elección de $c$ no afecta ni a la convergencia ni al valor. La definición exige que converjan las dos colas por separado: $\int_{-\infty}^{\infty}x\dif x$ diverge, aunque $\int_{-R}^{R}x\dif x=0$ para todo $R$.

\begin{proposicion}[Integrales $p$ en el infinito]\label{prop:p-infinito}
Sea $p\in\R$. La integral $\int_1^{\infty}\frac{\dif x}{x^{p}}$ converge si y solo si $p>1$, y en ese caso vale $\frac{1}{p-1}$.
\end{proposicion}

\addtocounter{definicion}{1}%

\section{Integrandos no acotados}

\begin{definicion}[Integral impropia de segunda especie]\label{def:impropia2}
Sea $f\colon(a,b]\to\R$ integrable en $[c,b]$ para todo $c\in(a,b)$ y no acotada cerca de $a$. Se define $\int_a^{b}f=\lim_{c\to a^{+}}\int_c^{b}f$ si el límite existe y es finito. Análogamente si la singularidad está en $b$. Si hay una singularidad en un punto interior, o varias, se divide el intervalo y se exige que converja cada trozo.
\end{definicion}

\begin{proposicion}[Integrales $p$ en el origen]\label{prop:p-cero}
La integral $\int_0^{1}\frac{\dif x}{x^{p}}$, con $p>0$, converge si y solo si $p<1$, y en ese caso vale $\frac{1}{1-p}$.
\end{proposicion}

Las \cref{prop:p-infinito,prop:p-cero} son simétricas: $x^{-p}$ es integrable en el infinito si decrece más deprisa que $1/x$, y en el origen si crece más despacio que $1/x$. En particular, $\int_0^{\infty}x^{-p}\dif x$ diverge para todo $p$.

\addtocounter{definicion}{1}%

\section{Criterios de convergencia}

Casi nunca se conoce una primitiva explícita, y conviene decidir la convergencia comparando con integrales conocidas. Los criterios se enuncian en $[a,\infty)$; para las de segunda especie son idénticos.

\begin{teorema}[Criterio de comparación]\label{teo:comparacion-int}
Sean $f,g$ integrables en $[a,b]$ para todo $b>a$, con $0\le f(x)\le g(x)$ para $x\ge a$. Si $\int_a^{\infty}g$ converge, $\int_a^{\infty}f$ converge y $\int_a^{\infty}f\le\int_a^{\infty}g$. Si $\int_a^{\infty}f$ diverge, $\int_a^{\infty}g$ diverge.
\end{teorema}

\begin{teorema}[Comparación en el límite]\label{teo:comparacion-limite}
Sean $f,g>0$ en $[a,\infty)$, integrables en cada $[a,b]$, con $\lim_{x\to\infty}\frac{f(x)}{g(x)}=L\in(0,\infty)$. Entonces $\int_a^{\infty}f$ y $\int_a^{\infty}g$ son ambas convergentes o ambas divergentes.
\end{teorema}

La demostración se pide en el \cref{eje:4-limite-comp}. En la práctica se compara con $x^{-p}$: por ejemplo, $\frac{1}{\sqrt{x^{3}+1}}$ se comporta como $x^{-3/2}$ en el infinito, y $\int_1^{\infty}\frac{\dif x}{\sqrt{x^{3}+1}}$ converge.

\begin{teorema}[Convergencia absoluta]\label{teo:absoluta-int}
Sea $f$ integrable en cada $[a,b]$. Si $\int_a^{\infty}\abs{f}$ converge, también converge $\int_a^{\infty}f$, y $\bigl|\int_a^{\infty}f\bigr|\le\int_a^{\infty}\abs{f}$.
\end{teorema}

\addtocounter{definicion}{1}%

\section{La función Gamma}

Para todo $n\in\N_0$, el límite $\lim_{x\to\infty}x^{n}e^{-x}=0$ (módulo~0.5) expresa que la exponencial domina a cualquier potencia. Gracias a ello las integrales $\int_0^{\infty}x^{n}e^{-x}\dif x$ convergen.

\begin{definicion}[Función Gamma]\label{def:gamma}
Para $s>0$ se define $\Gamma(s)=\int_0^{\infty}x^{s-1}e^{-x}\dif x$.
\end{definicion}

La integral es impropia en el infinito y, si $s<1$, también en el origen. Converge para todo $s>0$: en $(0,1]$, $0<x^{s-1}e^{-x}\le x^{s-1}=x^{-(1-s)}$ con $1-s<1$ (\cref{prop:p-cero}); en $[1,\infty)$, como $x^{s+1}e^{-x}\to0$, existe $x_0$ con $x^{s-1}e^{-x}\le x^{-2}$ para $x\ge x_0$ (\cref{prop:p-infinito}). El \cref{eje:4-gamma-conv} prueba que diverge si $s\le0$.

\begin{proposicion}\label{prop:gamma-factorial}
Para todo $s>0$, $\Gamma(s+1)=s\,\Gamma(s)$. En particular, $\Gamma(1)=1$ y $\Gamma(n+1)=n!$ para todo $n\in\N_0$.
\end{proposicion}

La función Gamma extiende el factorial a valores no enteros; por ejemplo, $\Gamma(\frac12)=\sqrt\pi$ (\cref{eje:4-gamma-medio}).

\section{La integral gaussiana}

La función $e^{-x^{2}}$ no tiene primitiva elemental, pero su integral sobre toda la recta puede calcularse exactamente. Es la constante que normaliza la distribución normal.

\begin{teorema}[Integral gaussiana]\label{teo:gauss}
$\displaystyle\int_{-\infty}^{\infty}e^{-x^{2}}\dif x=\sqrt\pi$.
\end{teorema}

\begin{observacion}[Sobre el paso a coordenadas polares]
La demostración usa dos hechos del cálculo en varias variables: que el producto de dos integrales es una integral doble (teorema de Fubini) y la fórmula del cambio a coordenadas polares, $\dif x\dif y=r\dif r\dif\theta$, de la que el argumento de las coronas es la versión intuitiva. Ambos se justifican en el módulo~0.7; véase también la §14.2 «Changing to Better Coordinates» de Strang~\cite{strang}. El \cref{eje:4-wallis-gauss} obtiene el mismo resultado sin salir de una variable, con las integrales de Wallis del \cref{eje:3-wallis}.
\end{observacion}

\begin{corolario}\label{cor:gauss}
Para $\mu\in\R$ y $\sigma>0$,
\[
\int_{-\infty}^{\infty}e^{-x^{2}/2}\dif x=\sqrt{2\pi},\qquad\int_{-\infty}^{\infty}e^{-(x-\mu)^{2}/(2\sigma^{2})}\dif x=\sigma\sqrt{2\pi}.
\]
En consecuencia, $\varphi(x)=\frac{1}{\sigma\sqrt{2\pi}}e^{-(x-\mu)^{2}/(2\sigma^{2})}$ es una densidad de probabilidad: es positiva e integra~1.
\end{corolario}

\addtocounter{definicion}{1}%

\addtocounter{definicion}{1}%

